\chapter{Partial observation, estimation and policy value}\label{ch:pomdp}
\section{A precise event clock}
Repeated bank decisions may depend on unobserved financial resilience,
liquidity needs or customer intent. A POMDP represents a latent state $S_t$,
an observation $O_t$ available before action $A_t$, a reward $R_t$, and
the next state and observation. Let
$H_t=(O_0,A_0,R_0,\ldots,O_t)$ be the recorded decision history.
The policy available to the bank must be measurable with respect to this
history, or to an explicitly enlarged logged information set.

For clarity, suppose a finite-state model has initial law $\mu_\theta$,
initial emission $B_\theta(o_0\mid s_0)$, transition
$P_\theta(s'\mid s,a)$, reward density $F_\theta(r\mid s,a,s')$
and next emission $B_\theta(o'\mid s',a)$. These factorization choices
are modeling assumptions. A joint emission for rewards and signals may be
needed when conditional independence is inappropriate.

Let the behavior policy be $q_\eta(a\mid s,h)$, potentially using private
information summarized by $s$. The observed-data likelihood for one
trajectory is
\begin{align}
 L(\theta,\eta)=\sum_{s_{0:T}}&
 \mu_\theta(s_0)B_\theta(o_0\mid s_0)\nonumber\\
 {}\times\prod_{t=0}^{T-1}&
 q_\eta(a_t\mid s_t,h_t)
 P_\theta(s_{t+1}\mid s_t,a_t)
 F_\theta(r_t\mid s_t,a_t,s_{t+1})
 B_\theta(o_{t+1}\mid s_{t+1},a_t).
 \label{eq:pomdplikelihood}
\end{align}
Omitting the behavior factor is justified for estimating environment
parameters when action assignment depends only on the recorded history
and is parameter-distinct: it can then be factored out of the latent-state
sum in the joint likelihood. If the logger uses unobserved information,
the action itself is evidence about latent state and the factor cannot
generally be removed. Fitting a model conditional on the action sequence
without this distinction can misinterpret selection as a transition effect.

\section{Filtering, smoothing and estimation}
The belief-state construction connects partial observation to planning
\citep{kaelbling1998}; estimation additionally requires an explicit likelihood.
For an observed-history policy, after observing $O_t$ and previous rewards,
let $b_t(s)=\Prob_\theta(S_t=s\mid H_t)$.
The action is not an additional latent-state signal conditional on $H_t$.
After observing $r_t,o_{t+1}$,
\begin{equation}
 b_{t+1}(s')=
 \frac{B_\theta(o_{t+1}\mid s',a_t)
 \sum_s b_t(s)P_\theta(s'\mid s,a_t)F_\theta(r_t\mid s,a_t,s')}
 {\sum_{u}B_\theta(o_{t+1}\mid u,a_t)
 \sum_s b_t(s)P_\theta(u\mid s,a_t)F_\theta(r_t\mid s,a_t,u)}.
 \label{eq:belief}
\end{equation}
If the logger depends on latent state, first update
$b_t^+(s)\propto q_\eta(a_t\mid s,H_t)b_t(s)$ and use $b_t^+$ in
the observation update. Under a new intervention policy based only on
recorded history, this observational action update is removed.

Normalizers from the forward recursion accumulate the log likelihood,
avoiding numerical underflow. Backward messages produce smoothed state
probabilities and expected transition counts. In an EM algorithm,
\[
 \widehat P_a(s,s')=
 \frac{\sum_{i,t:A_{it}=a}\xi_{it}(s,s')}
      {\sum_{i,t:A_{it}=a}\sum_u\xi_{it}(s,u)}
\]
updates tabular transitions, with explicit prior/regularization when
counts are small. Emission and reward parameters are updated consistently
with their likelihood. Unobserved-logger parameters, if included, require
their own model and identification restrictions.

Alternatives include direct differentiation of the marginalized likelihood,
Bayesian inference, variational approximations and particle methods for
nonlinear continuous states. Compare multiple initializations, held-out
predictive likelihood, residual dynamics, posterior predictive checks,
parameter recovery and prior sensitivity. Label permutations require a
normalization or permutation-invariant comparison. Numerical convergence
and good predictive fit do not establish causal identification.

\section{Statistical and causal identification differ}
\citet{hmmident2016} establish HMM identification under full-rank ergodic
transitions, distinct emissions and other stated conditions, up to state
labels. This is a result about an observed distribution within an HMM
class. A controlled bank model additionally needs sufficiently informative
actions, identifiable reward/observation mechanisms and an intervention
interpretation. The stationary uncontrolled theorem does not supply all
of these conditions.

\begin{proposition}[Observed fit need not identify an action effect]
There exist two binary-action binary-outcome models with the same strictly
positive observed law and different average treatment effects.
\end{proposition}
\begin{proof}
In model I, $U\sim\mathrm{Bernoulli}(1/2)$,
$\Prob(A=1\mid U)=0.2+0.6U$ and
$\Prob(Y=1\mid U,A)=0.2+0.6U$. Action has no causal effect.
In model II, $A\sim\mathrm{Bernoulli}(1/2)$ independently and
$\Prob(Y=1\mid A)=0.32+0.36A$. Its effect is $0.36$.
Both imply
$\Prob(A,Y)=(0.34,0.16,0.16,0.34)$ in cell order
$(0,0),(0,1),(1,0),(1,1)$. Their observational likelihoods agree
while their intervention distributions differ.
\end{proof}

This one-step counterexample embeds in a sequential problem. Observed
action positivity and perfect predictive fit therefore cannot guarantee
causal policy identification. The practical remedy is an explicit route:
randomized or measured-history assignment, defensible structural restrictions,
valid instruments/proxies, or partial identification.

\section{Proxies and instruments are measurement designs}
In a static proximal model let $U$ be hidden confounding, $Z$ a
treatment-inducing proxy, $W$ an outcome-inducing proxy and $X$ observed
covariates. One bridge equation is
\begin{equation}
 \E[Y\mid Z,A,X]=\E[h(W,A,X)\mid Z,A,X].
 \label{eq:bridge}
\end{equation}
Under latent exchangeability, proxy exclusions such as
$Y\perp Z\mid A,U,X$ and $W\perp(A,Z)\mid U,X$, completeness,
bridge existence and support, a causal mean can be expressed as
$\E[h(W,a,X)]$ \citep{miao2018,proximal2020}.
The exclusions describe how variables are generated. Two highly predictive
features are not necessarily two valid proxies.

The bridge is an inverse problem. Weak proxies can make it ill-conditioned;
regularization trades variance against bias. Minimax methods and regression
bridges provide estimation strategies under corresponding restrictions
\citep{ghassami2022,proxreg2024}. The product should show proxy-role
justification, bridge residuals, conditioning diagnostics and regularization
sensitivity; these do not empirically prove completeness.

POMDP/proximal OPE theories impose particular information and proxy
structures \citep{tennenholtz2020,bennett2021,bennettprox2023}.
IV-based sequential evaluation also requires relevance, exclusion,
interaction restrictions and dynamic conditional independences
\citep{xu2023}. These are research-grade options for a bank with the
right measurement design, not generic repairs for absent action logs.
Recent hidden-treatment work similarly requires informative surrogates
and additional proxy restrictions \citep{hidden2026}.

\section{Policy evaluation before optimization}
The actor must also be clear. A POMDP policy can describe a bank choosing
actions given customer-state uncertainty. A customer's refinancing or
account-closing choice is a different behavior model, potentially with
private preferences, switching costs and beliefs. Estimating that choice
model requires restrictions on utility, expectations and state transitions.
Optimizing the bank's supplied reward function does not identify customer
preferences. Structural demand or dynamic choice components can supply
this behavior only with their own identifying evidence.

With sufficient observed histories and unconfounded logging, a target
policy $\pi$ has finite-horizon value
\[
 V^\pi=\E^\pi\!\left[\sum_{t=0}^{T-1}\gamma^tR_t\right].
\]
Importance ratios $\rho_t=\pi(A_t\mid H_t)/e(A_t\mid H_t)$
and their products correct behavior-policy selection when the relevant
conditions hold. A sequential DR value score can combine action-value
regression with weighted Bellman residuals. For a Markov state $X_t$,
one familiar finite-horizon form is
\begin{equation}
 \widehat V^\pi=
 \frac1n\sum_i\left[
 \widehat V_0(X_{i0})+
 \sum_{t=0}^{T-1}\gamma^tW_{it}
 \{R_{it}+\gamma\widehat V_{t+1}(X_{i,t+1})
                   -\widehat Q_t(X_{it},A_{it})\}\right],
 \label{eq:opedr}
\end{equation}
where $\widehat V_t(x)=\sum_a\pi_t(a\mid x)\widehat Q_t(x,a)$,
$\widehat V_T=0$, and $W_{it}=\prod_{s\leq t}\rho_{is}$.
Markov restrictions permit more efficient marginalized-ratio approaches
\citep{kallus2020}; replacing $X_t$ with an estimated belief does not
automatically inherit their identification or variance theory.

For two fixed policies, estimate the difference with paired scores and
preserve their covariance. A candidate deployment rule is a positive
lower confidence bound on incremental net value, with customer-outcome
constraints and correction for policy selection. For learned policies,
use held-out trajectories or valid uniform inference; ordinary pointwise
intervals after a large policy search are optimistic
\citep{policy2021}.

An ambiguity set $\mathcal M$ gives a conservative criterion
$\inf_{M\in\mathcal M}\{V_M^\pi-V_M^{\pi_0}\}$.
Its meaning depends on why $\mathcal M$ contains the relevant model.
Sharp dynamic bounds in \citet{sharpdynamic2026} assume stage-local
confounding; persistent customer fragility can violate the rectangular
structure required for local backward optimization. A recent negative result
constructs exponentially difficult OPE despite favorable behavior-marginal
coverage under history-dependent logging \citep{opehard2026}. It cautions
against diagnosing difficulty with only one-step action counts.

\section{A reproducible estimation and filtering illustration}
The supplied two-state model has initial law $(0.7,0.3)$, emissions
\[
 B=\begin{pmatrix}.8&.2\\.2&.8\end{pmatrix},\quad
 P_0=\begin{pmatrix}.9&.1\\.25&.75\end{pmatrix},\quad
 P_1=\begin{pmatrix}.95&.05\\.55&.45\end{pmatrix}.
\]
For actions $(1,0,1)$ and observations $(1,0,1,1)$, forward
likelihood $0.0234985$ equals explicit enumeration of all 16 latent
paths. Rewards are omitted from filtering in this illustration.
Expected immediate reward vectors are $(12,-8)$ for no assistance and
$(10,1)$ for assistance. With fragile-state belief $b$, assistance
is myopically preferable when $-2+11b>0$, or $b>2/11$.
This threshold is not the long-horizon optimal policy.

\begin{figure}[htbp]\centering
\includegraphics[width=\textwidth]{figures/belief_value.pdf}
\caption{Exact filtering and a myopic value comparison in a supplied causal model.}
\label{fig:belief}
\end{figure}

A separate simulation generates 700 trajectories of 14 randomized actions.
With emissions and initial law fixed and known, likelihood optimization
recovers the four transition-to-fragile probabilities as approximately
$(0.0935,0.7002,0.0774,0.4192)$ against truth
$(0.1,0.75,0.05,0.45)$. Three starts converge to closely agreeing
objectives; maximum absolute error is $\NumHMMError$.
This verifies a small controlled estimation construction. It does not
establish unknown-emission identification, uncertainty coverage, real-bank
state recovery or bank policy value.

Counterfactual trajectories introduce another distinction.
A Gumbel-max structural model specifies an exogenous-noise coupling across
actions \citep{oberst2019}. Different couplings can share transition
probabilities while producing different individual counterfactual paths.
The product should label such paths as model-conditional scenarios.
